\documentclass[12pt,a4paper]{article}

% ---------- Packages ----------
\usepackage[utf8]{inputenc}
\usepackage[T5]{fontenc}
\usepackage[vietnamese]{babel}
\usepackage{lmodern}
\usepackage[margin=2.2cm]{geometry}
\usepackage{graphicx}
\usepackage{float}
\usepackage{tikz}
\usetikzlibrary{calc}
\usepackage{booktabs}
\usepackage{longtable}
\usepackage{tabularx}
\usepackage{array}
\usepackage[table]{xcolor}
\usepackage{listings}
\usepackage{enumitem}
\usepackage{fancyhdr}
\usepackage{titlesec}
\usepackage[hidelinks]{hyperref}

\definecolor{bkblue}{RGB}{3,43,145}
\definecolor{codebg}{RGB}{246,248,250}
\newcommand{\todo}[1]{\textcolor{red}{\textbf{[#1]}}}
\newcolumntype{L}[1]{>{\raggedright\arraybackslash}p{#1}}

\lstset{
  language=Java, basicstyle=\ttfamily\small, backgroundcolor=\color{codebg},
  frame=single, rulecolor=\color{black!20}, numbers=left,
  numberstyle=\tiny\color{gray}, keywordstyle=\color{bkblue}\bfseries,
  commentstyle=\color{green!45!black}\itshape, stringstyle=\color{red!60!black},
  breaklines=true, showstringspaces=false, columns=fullflexible, xleftmargin=1.5em
}

\titleformat{\section}{\Large\bfseries\color{bkblue}}{\thesection.}{0.6em}{}
\titleformat{\subsection}{\large\bfseries}{\thesubsection.}{0.6em}{}

\pagestyle{fancy}
\fancyhf{}
\lhead{Kiểm thử phần mềm -- Nhóm HEHE2027}
\rhead{Bài tập lớn 1: Code Review}
\cfoot{\thepage}
\setlength{\headheight}{15pt}

\sloppy
\setlength{\emergencystretch}{3em}

\begin{document}

% ================= COVER =================
\begin{titlepage}
\begin{tikzpicture}[overlay,remember picture]
\draw [line width=2.25pt]
    ($ (current page.north west) + (2.0cm,-1.5cm) $)
    rectangle
    ($ (current page.south east) + (-1.5cm,2cm) $);
\draw [line width=0.5pt]
    ($ (current page.north west) + (2.2cm,-1.7cm) $)
    rectangle
    ($ (current page.south east) + (-1.7cm,2.2cm) $);
\end{tikzpicture}
\begin{center}
\textbf{ĐẠI HỌC QUỐC GIA THÀNH PHỐ HỒ CHÍ MINH} \\
\textbf{TRƯỜNG ĐẠI HỌC BÁCH KHOA}\\
\textbf{KHOA KHOA HỌC VÀ KỸ THUẬT MÁY TÍNH}

\vspace{1cm}
\includegraphics[width=0.32\textwidth]{img/logo/logo.png}
\vspace{1cm}

\begin{tabular}{c}
    \textbf{\Large KIỂM THỬ PHẦN MỀM (SOFTWARE TESTING)}\\[0.4cm]
    \textbf{\large BÀI TẬP LỚN 1: SOFTWARE TECHNICAL REVIEW}\\[0.2cm]
    \large Code review dự án \textit{BirdBook}\\[0.4cm]
    \textbf{Học kỳ 1, năm học 2026--2027 -- Lớp L01 -- Nhóm HEHE2027}\\
    \hline
\end{tabular}

\vspace{1cm}
\begin{tabular}{p{3cm} p{6cm} l}
    \textbf{GVHD:} & Trần Thị Ngọc Trâm & \\[0.2cm]
    \textbf{Sinh viên:} & Nguyễn Đức Phương & 2312749 \\
                        & Lê Sĩ Nhật Quang & 2312791 \\
                        & Huỳnh Đức Tài & 2212971 \\
                        & Đào Xuân Sơn & 2312954 \\
                        & Sơn Văn Phát & 2212535 \\
\end{tabular}

\vfill
\fontsize{14pt}{0pt}\selectfont TP. HỒ CHÍ MINH, THÁNG 9/2026
\end{center}
\end{titlepage}

\tableofcontents
\newpage

% ================= 1 =================
\section{Giới thiệu}
Nhóm thực hiện review kỹ thuật (\textit{technical review}) mã nguồn Java của \textit{BirdBook},
một hệ thống đặt vé máy bay gồm \textbf{server} (Spring Boot: MVC, JPA, Security, Cache với Redis)
và \textbf{client} (JavaFX). Việc review dựa trên \textit{Java Code Review Checklist} gồm 43 mục,
chia thành ba nhóm: \textit{Defect} (mục 1--38), \textit{Ambiguity} (39--40) và \textit{Redundancy} (41--43).

Tài liệu nộp kèm là file \texttt{HEHE2027\_L01\_Proj\#1-Code Review.xlsx}, gồm sheet \texttt{CoverPage},
mỗi file một sheet review và sheet \texttt{All problems} tổng hợp tất cả lỗi.

% ================= 2 =================
\section{Phân công công việc}
Mỗi file có một người review (1st-reviewer) và một người kiểm tra lại (2nd-reviewer).
Mỗi thành viên review tối thiểu 3 file.

{\small
\begin{longtable}{|c|L{3.1cm}|L{5.4cm}|c|c|L{2.2cm}|}
\hline
\rowcolor{gray!15}
\textbf{STT} & \textbf{Thành viên} & \textbf{1st-reviewer (file)} & \textbf{LoC} & \textbf{Lỗi} & \textbf{Đóng góp} \\ \hline
\endhead
1 & Nguyễn Đức Phương \newline 2312749 &
  \texttt{VilleServiceImpl}, \texttt{AdminStatsController}, \texttt{UserDetailsImpl} & 329 & 28 & 20\% \\ \hline
2 & Lê Sĩ Nhật Quang \newline 2312791 &
  \texttt{AdminBilletController}, \texttt{AdminClientController}, \texttt{AdminVilleController} & 288 & 26 & 20\% \\ \hline
3 & Huỳnh Đức Tài \newline 2212971 &
  \texttt{AuthController} (server), \texttt{CategorieController}, \texttt{ReservationController}, \texttt{SearchController} & 308 & 29 & 20\% \\ \hline
4 & Đào Xuân Sơn \newline 2312954 &
  \texttt{BookingDetailsController}, \texttt{SignupController}, \texttt{SideBarController} & 325 & 31 & 20\% \\ \hline
5 & Sơn Văn Phát \newline 2212535 &
  \texttt{PersonalInfoController}, \texttt{PassengerEditorController}, \texttt{SearchResultsController} & 342 & 31 & 20\% \\ \hline
\multicolumn{3}{|r|}{\textbf{Tổng: 16 file}} & \textbf{1592} & \textbf{145} & \textbf{100\%} \\ \hline
\end{longtable}
}

\noindent\textbf{Phân công 2nd-reviewer (validate):}
\begin{itemize}[noitemsep]
  \item Nguyễn Đức Phương validate các file của Lê Sĩ Nhật Quang.
  \item \todo{Bổ sung người validate cho các file của Phương, Tài, Sơn, Phát}
\end{itemize}

% ================= 3 =================
\section{Quy trình review}
\begin{enumerate}[noitemsep]
  \item \textbf{Chuẩn bị:} mỗi thành viên đọc file được giao cùng các interface, entity, DTO và service
        liên quan để hiểu ngữ cảnh.
  \item \textbf{Review theo checklist:} với từng mục trong 43 mục, ghi \texttt{Yes}/\texttt{No}/\texttt{\#NA} vào cột
        \textit{YES/NO/\#N/A}. Cột \textit{Objective} chứa đáp án mong đợi; cột \textit{Pass?} tự so khớp hai cột
        và ra \texttt{No} khi mục đó có lỗi. Với mỗi lỗi, ghi số dòng, mô tả và đề xuất sửa.
        \texttt{\#NA} dùng cho mục không áp dụng (ví dụ file không có vòng lặp hay \texttt{switch}).
  \item \textbf{Kiểm chứng:} 2nd-reviewer đối chiếu từng lỗi với mã nguồn để loại lỗi ``giả''.
  \item \textbf{Tổng hợp:} chép lỗi sang sheet \texttt{All problems}; các lỗi nghiêm trọng nằm ngoài checklist
        được đưa vào Mục~5 (Extra).
\end{enumerate}

% ================= 4 =================
\section{Kết quả review}

\subsection{Tổng quan}
Nhóm review \textbf{16 file (1592 LoC)} và ghi nhận \textbf{145 lỗi}. Các mục checklist bị vi phạm nhiều nhất:

\begin{table}[H]
\centering
\small
\begin{tabular}{|c|L{8.8cm}|c|}
\hline
\rowcolor{gray!15}\textbf{Mục} & \textbf{Nội dung} & \textbf{Số file vi phạm} \\ \hline
28, 30, 32 & Thiếu header comment / comment cho biến / comment nói chung & 16 / 16 / 16 \\ \hline
26 & Xử lý exception không phù hợp & 10 \\ \hline
34 & Định dạng, thụt lề không nhất quán & 10 \\ \hline
1, 40 & Tên biến không rõ nghĩa, sai quy ước & 9 / 8 \\ \hline
41, 43 & Biến thừa / code lặp lại & 8 / 8 \\ \hline
4 & Giá trị nên là hằng số & 7 \\ \hline
36 & Tính toán lại giá trị nhiều lần & 6 \\ \hline
\end{tabular}
\caption{Các mục checklist vi phạm nhiều nhất}
\end{table}

Nhận xét chung: hầu như toàn bộ mã nguồn \textbf{không có comment}; exception thường bị ``nuốt''
(\texttt{catch} rỗng hoặc chỉ \texttt{printStackTrace()}); tên biến viết tắt (\texttt{vr}, \texttt{clt}, \texttt{ss}, \texttt{vd}).
Ngoài ra có một số lỗi làm chương trình \textbf{không biên dịch được} hoặc \textbf{chạy sai}, trình bày dưới đây.

\subsection{Các lỗi nghiêm trọng tiêu biểu}

{\footnotesize
\begin{longtable}{|L{3.1cm}|c|L{6.3cm}|L{4.0cm}|}
\hline
\rowcolor{gray!15}\textbf{File (dòng)} & \textbf{Mục} & \textbf{Lỗi} & \textbf{Đề xuất} \\ \hline
\endhead
\texttt{PassengerEditor\-Controller} (49) & 16 & \texttt{if (mode = Mode.CREATE)} là phép gán, không phải so sánh, nên file \textbf{không biên dịch được}. & \texttt{mode == Mode.CREATE} \\ \hline
\texttt{AdminVille\-Controller} (68, 75) & 39 & Tham số tên \texttt{ID} nhưng code dùng \texttt{id}, nên \textbf{lỗi biên dịch}. & Thống nhất dùng \texttt{id}. \\ \hline
\texttt{AdminStats\-Controller} (87--90) & 15, 13 & \texttt{100 / totalSeats} với \texttt{totalSeats = 0} gây \texttt{ArithmeticException} ở \textbf{mọi} request \texttt{/plane-occupancy}; phép chia \texttt{int/int} còn làm mất phần thập phân. & Xoá code thử; lấy tỉ lệ từ \texttt{PlaneOccupancyStatsDTO}. \\ \hline
\texttt{AdminStats\-Controller} (124) & 17 & \texttt{year!=null \&\& threshold!=null || quarter!=null}: chỉ cần có \texttt{quarter} là service bị gọi với \texttt{year = null}. & \texttt{year != null \&\& threshold != null} \\ \hline
\texttt{AuthController} server (52) & 17 & \texttt{authentication == null \&} \texttt{authentication.\allowbreak isAuthenticated()}: toán tử \texttt{\&} vẫn tính vế phải, nên gây \texttt{NullPointerException} khi \texttt{authentication} là \texttt{null}. & Dùng toán tử ngắn mạch \texttt{\&\&}/\texttt{||} và sửa lại điều kiện. \\ \hline
\texttt{AuthController} server (72) & 25 & Method có kiểu trả về \texttt{ResponseEntity<ClientDTO>} nhưng thiếu lệnh \texttt{return}. & Trả về kết quả của \texttt{register()}. \\ \hline
\texttt{Reservation\-Controller} (45) & 17 & Dùng \texttt{|} thay vì \texttt{||}. & Dùng \texttt{||}. \\ \hline
\texttt{PersonalInfo\-Controller} (68--69) & 27 & \texttt{setPrenom(getNom())} gán họ vào tên; \texttt{setPassword(fName.getText())} ghi đè mật khẩu người dùng đã nhập. & Gán tên bằng \texttt{fName.\allowbreak getText()}; bỏ lệnh \texttt{setPassword}. \\ \hline
\texttt{PersonalInfo\-}, \texttt{PassengerEditor\-Controller} & 26 & \texttt{DatePicker.getValue().toString()} khi chưa chọn ngày gây \texttt{NullPointerException}. & Kiểm tra ngày trước khi gửi. \\ \hline
\texttt{SideBar\-Controller} (48--61) & 11 & Gọi \texttt{searches.get(0..2)} mà không kiểm tra kích thước danh sách; lỗi bị \texttt{catch} rỗng che đi. & Kiểm tra \texttt{size()} trước khi truy cập. \\ \hline
\texttt{SearchResults\-Controller} (27, 115--123) & 41, 27 & \texttt{departureDate} được gán nhưng không bao giờ dùng, nên kết quả không lọc theo ngày; \texttt{loadFlightInfo(null)} truyền id \texttt{null}. & Áp dụng ngày vào truy vấn; truyền id chuyến bay đã chọn. \\ \hline
\texttt{VilleServiceImpl} (86--93) & 26 & \texttt{findById} trả về \texttt{null} thay vì ném \texttt{NotFoundException}; kiểm tra \texttt{Optional != null} luôn đúng. & \texttt{orElseThrow(...)} \\ \hline
\end{longtable}
}

% ================= 5 =================
\section{Extra: lỗi nghiêm trọng ngoài checklist}

\paragraph{E1. Cache không hoạt động do gọi method trong cùng class (VilleServiceImpl, dòng 98, 107, 115).}
Spring xử lý \texttt{@CachePut}/\texttt{@CacheEvict} thông qua proxy, nên lời gọi \texttt{this.deleteById()} hay
\texttt{this.updateDTO()} từ bên trong class sẽ \textbf{bỏ qua cache}. Hậu quả: sau \texttt{delete(ville)},
\texttt{findDTOById(id)} vẫn trả về thành phố đã bị xoá. \textit{Đề xuất:} đặt annotation cache trực tiếp trên
\texttt{delete()}/\texttt{update()}, hoặc gọi qua bean được inject.

\paragraph{E2. Mã nguồn không biên dịch được.}
Khi biên dịch thử phần server, nhóm phát hiện: \texttt{AdminPassagerController.java} thiếu dấu \texttt{\}} cuối file;
một số file import sai tên package \texttt{jakarta.valIDation} và \texttt{org.springframework.valIDation} (đúng phải là
\texttt{validation}, có vẻ do thay thế hàng loạt ``id'' thành ``ID''). Cùng với hai lỗi biên dịch ở Mục 4.2, dự án hiện
\textbf{không build được}. \textit{Đề xuất:} bật build trên CI (Maven/Gradle) để chặn các commit không biên dịch được.

\paragraph{E3. Không kiểm tra dữ liệu đầu vào (AdminStatsController).}
\texttt{quarter}, \texttt{threshold} và \texttt{year} không được kiểm tra khoảng giá trị (quý 1--4, ngưỡng 0--100\%).
\textit{Đề xuất:} dùng Bean Validation (\texttt{@Min}, \texttt{@Max}, \texttt{@DecimalMin}, \texttt{@DecimalMax}) với \texttt{@Validated}.

\paragraph{E4. Trạng thái \texttt{static} dùng chung giữa các màn hình (client).}
\texttt{BookingDetailsController}, \texttt{SignupController} và \texttt{PassengerEditorController} truyền dữ liệu giữa
các màn hình qua field \texttt{static} (có chỗ còn là \texttt{public static}), nên dễ bị dữ liệu cũ hoặc lẫn dữ liệu khi
mở nhiều cửa sổ. \textit{Đề xuất:} truyền dữ liệu qua instance controller (\texttt{FXMLLoader.getController()}).

% ================= 6 =================
\section{Đề xuất chung}
\begin{enumerate}[noitemsep]
  \item Bắt buộc code phải \textbf{build thành công} và qua code review trước khi merge; xoá code thử (\texttt{//just try}, \texttt{System.out.println}).
  \item Thống nhất cách xử lý lỗi: service ném exception, \texttt{GlobalExceptionHandler} xử lý; client hiển thị thông báo thay vì \texttt{catch} rỗng.
  \item Validate dữ liệu đầu vào ở cả client (form) và server (Bean Validation).
  \item Dùng công cụ phân tích tĩnh (SonarLint, SpotBugs, Checkstyle) để tự động bắt lỗi quy ước đặt tên, biến thừa, chia cho 0, dùng \texttt{\&}/\texttt{|} thay cho \texttt{\&\&}/\texttt{||}.
  \item Bổ sung Javadoc cho các lớp public và thống nhất định dạng code (formatter chung của IDE).
\end{enumerate}

% ================= 7 =================
\section{AI Usage Declaration}

\subsection{Nguyễn Đức Phương}
\textbf{Công cụ:} Claude (Anthropic), dùng qua ứng dụng Claude Code (model Claude Opus 5.5).

\noindent\textbf{Prompt đã sử dụng} (nguyên văn):
\begin{enumerate}[noitemsep]
  \item ``hãy mô tả lại đề bài''
  \item ``giải thích ý nghĩa của 3 file này''
  \item ``giải thích các lỗi ở từng file, tại sao lại lỗi''
  \item ``ở file đầu tiên mục comment thì tất cả phải là no hết chứ nhỉ, vì đấu có cmt nào đâu''
  \item ``ở sheet all problem, hãy copy các giá trị từ 3 file đã review qua tương ứng với ví dụ cho toi''
  \item ``liệt kê những thay đổi của từng file''
\end{enumerate}

\noindent\textbf{Kết quả, nhận xét và đề xuất của AI:}
\begin{itemize}[noitemsep]
  \item Tóm tắt đề bài.
  \item Review 3 file, đưa ra 28 lỗi theo checklist cùng các lỗi Extra E1 và E3, kèm đề xuất sửa; giải thích nguyên nhân từng lỗi.
  \item Viết script Python điền kết quả vào template Excel; viết phiên bản code đã sửa của 3 file.
\end{itemize}

\subsection{Các thành viên khác}
\todo{Lê Sĩ Nhật Quang, Huỳnh Đức Tài, Đào Xuân Sơn, Sơn Văn Phát: mỗi người ghi công cụ AI, prompt và kết quả;
nếu không dùng AI, ghi: ``We hereby confirm that no AI assistance was used in this work.''}

\end{document}
